\documentclass[a4paper,11pt]{article} % Don't change this
\usepackage{fullpage} % Don't change this

\usepackage{titling} % Don't change this
\pretitle{\vspace{-6em}\begin{center}\larger[2]\bf} % Don't change this
\postauthor{\end{center}} % Don't change this
\preauthor{\vspace{-2em}\begin{center}\smaller[1]} % Don't change this
\postauthor{\end{center}\vspace{-5em}} % Don't change this

\usepackage[shortlabels]{enumitem} % Don't change this
\setlist{itemsep=0em,topsep=0.3em} % Don't change this

%%% Some useful packages.
\usepackage{xcolor}
\usepackage{float}
\usepackage{amsmath,amssymb,mathtools}
\usepackage{graphicx}
\usepackage{subcaption}
\usepackage[hidelinks]{hyperref}
\usepackage{relsize}
\usepackage{cleveref}
\crefname{figure}{Fig.}{Figs.}

\usepackage[backend=bibtex,natbib=true,style=authoryear-icomp,maxnames=3,giveninits=true]{biblatex}
\citetrackerfalse
\bibliography{ref} % put bib entries in ref.bib

\usepackage[linesnumbered]{algorithm2e}
\usepackage{cleveref}
\crefname{equation}{Equation.}{Equations.}
\crefname{figure}{Figure}{Figures.}

\usepackage{lipsum}
\usepackage{wrapfig}
\usepackage{tikz}

%%% Add additional packages or macros below
%
% Everything above this line is proposal_template.tex verbatim. The course
% report template has the same preamble.
%
% How to use this file. Upload the folder report_v2/overleaf to Overleaf as it
% is, or build it locally with "tectonic main.tex" from this folder. It needs no
% file outside this folder. Every figure is a red placeholder box until the
% scripts of report_v2/EXPERIMENTS_TODO.md write PDFs into figures/.
%
% Red text is for the group and must be deleted before submission:
%   \TODO{...}        a task.
%   \pending{v}{ID}   a planning value from a read-only check on data/. Item ID
%                     of EXPERIMENTS_TODO.md recomputes it before it is final.
%   \budget{...}      the page budget of a section.
%   draftnote         a red block. The draft notes after the consent page hold
%                     the criteria self-assessment, the changelog and the
%                     reference checks.
\usepackage{booktabs}
\usepackage{longtable}
\usepackage{array}
\newcolumntype{L}[1]{>{\raggedright\arraybackslash}p{#1}}
\crefname{table}{Table}{Tables}
\newcommand{\file}[1]{\nolinkurl{#1}}

\newcommand{\TODO}[1]{\textcolor{red}{\textbf{[TODO: #1]}}}
\newcommand{\pending}[2]{\textcolor{red}{#1\textsuperscript{\,[#2]}}}
\newcommand{\budget}[1]{{\color{red}\footnotesize\textit{Page budget: #1.}\par}}
\newcommand{\placeholderfigure}[2][4.5cm]{%
  \fbox{\parbox[c][#1][c]{0.95\linewidth}{\centering\color{red}\small\textbf{Figure to make.} #2}}}
\newenvironment{draftnote}{\par\color{red}}{\par}
% A draft of what a section contributes, written as if the preliminary findings
% hold. Each block says why the finding is new and why it matters. Fold the
% text into the section's prose once the full runs confirm it, then delete the
% block.
\newenvironment{draftcontrib}[1]{\par\smallskip\noindent\color{red}%
  \textbf{Draft contribution, #1, if the preliminary findings hold.}\ }%
  {\ \textit{Verify with the full-run experiments before keeping any of this. If
  a finding it rests on does not hold, rewrite or drop the claim.}\par\smallskip}

\newcommand{\etalam}{\eta\lambda}
\newcommand{\tstar}{t^{\ast}}
\newcommand{\Fnorm}[1]{\lVert #1\rVert_F}

\begin{document}

\title{Scale or Shape? Decomposing NTK Movement During Grokking}
\author{Group G2 \\
Zain Al-Saffi (s4800836) $\cdot$
Jasper Chong (s4883420) $\cdot$
Aaron D'Mello (s4890257) \\
Zac Kienzle (s4803349) $\cdot$ Samuel Rieger (s4904959)}

\date{}
\maketitle

\begin{draftnote}
\noindent\textbf{Notes for the group. Delete this block before submission.}
\begin{itemize}
  \item \TODO{Later, replace every mention of ``rotation'' with ``shape'',
    including $R_t$'s name, the figure labels and the tier titles. The word
    ``rotation'' is kept for now so that the text matches what the group
    already uses.}
  \item This is a skeleton. Each section opens with the claim it has to
    support. Red numbers with an item ID are planning values from read-only
    checks on \file{data/} (5 October 2026). Recompute every one through its
    item in \file{EXPERIMENTS_TODO.md} before it becomes final.
  \item The main text must end by page 12. References and appendices do not
    count, and the appendices may take at most 5 pages. Each section shows its
    page budget.
  \item The draft notes after the consent page hold the criteria
    self-assessment, the changelog of what was cut and why, and the reference
    checks.
  \item Each section has a red block headed ``Draft contribution''. It says
    what the section adds and why it matters, written as if the preliminary
    findings hold. The blocks carry the narrative and answer the marking
    criteria for significance and novelty. Check each against the full-run
    results, fold what survives into the section's prose, and delete the block.
    Until then the blocks push the main text past its page budget.
\end{itemize}
\end{draftnote}

\begin{abstract}
  \TODO{Write last, from the finished sections. Keep one sentence for each
    point below.}
  Grokking is delayed generalisation, and what sets the length of the delay is
  disputed. Recent accounts time it by the weight norm, by the scale of the
  logits, or by the weight-decay rate. We split the movement of the neural
  tangent kernel into a change of scale and a rotation, and measure both, with
  the kernel's alignment to the labels, on 720 runs of a two-layer network
  learning modular addition. The runs cross laziness, width and weight decay.
  At a fixed width, grokking happens at a nearly fixed amount of rotation, but
  no statistic takes a fixed value across widths. The grokking time depends on
  both width and weight decay, which neither a rotation clock nor a
  weight-decay clock predicts alone. Read before any run groks, rotation is the
  best single predictor of the grokking time, and scale adds information that
  rotation lacks. \TODO{Final sentence: what this means for the accounts of
  grokking.}
\end{abstract}

% =====================================================================
\section{Introduction}
\budget{1.0 page}

Grokking is the delayed generalisation that
\citet{powerGrokkingGeneralizationOverfitting2022} found on small algorithmic
datasets. A network fits its training data early, and its test accuracy rises
from chance long afterwards. Modular addition learned by a two-layer perceptron
is a small instance of it. Full-batch gradient descent on the squared error
groks there with no explicit regularisation
\citep{gromovGrokkingModularArithmetic2023}. What sets the length of the delay
is an open question.

One account holds that grokking is a transition from lazy training to rich
training \citep{kumarGrokkingTransitionLazy2024}. In lazy training the network
behaves as a kernel method with its initial tangent kernel. In rich training the
kernel adapts to the task. The evidence that the kernel changes is a distance:
the empirical tangent kernel stays close to its initial value until the training
data are fitted, and moves far from it afterwards
\citep{mohamadiWhyYouGrok2024}. A distance does not say how the kernel changed.
A kernel can grow in scale while its eigenvectors stay fixed. It can also
rotate, so that its eigenvectors change.

The difference matters for the accounts of the grokking time. Recent work times
grokking by the weight norm \citep{khanhWeightNormSets2026}, by the scale of the
logits \citep{khanhWhatDoesWeight2026}, or by the weight-decay rate
\citep{kimGrokkingWeightDecayClock2026,pracherSpectralTheoryGrokking2026}. The
result that \citet{kumarGrokkingTransitionLazy2024} cite for weight decay forcing
the kernel to change is that of \citet{lewkowyczTrainingDynamicsDeep2021}, in
which the kernel only shrinks. None of these accounts measures whether the
kernel rotates.

We ask what sets the grokking time: the shape of the kernel, its scale, or the
weight-decay rate. Our hypothesis is that rotation sets it. We measure three
statistics of the empirical tangent kernel along training: its scale $S_t$, its
rotation $R_t$ and its alignment $A_t$ with the labels. We do so on a grid of
720 runs that crosses laziness, width and weight decay. Our contributions are as
follows.
\begin{enumerate}[1.]
  \item We split the movement of the kernel into a change of scale and a
    rotation, and use each one as a timing variable. On the baseline of
    \citet{kumarGrokkingTransitionLazy2024}, the usual distance is
    \pending{mostly a change of scale}{N2}.
  \item We test whether grokking happens at a fixed value of each statistic. A
    placebo test reads each run at another run's grokking step, so a statistic
    that merely stays in a narrow band does not pass.
  \item We cross width with weight decay and fit how the grokking time scales
    with each. A rotation clock and a weight-decay clock predict different
    exponents.
  \item We test whether the statistics, read before any run groks, predict when
    it will. Unlike the settings that produced them, these statistics can be
    measured in any network.
  \item \TODO{A theory contribution, if \cref{sec:tier4} keeps one. Its form
    depends on the option the group chooses there.}
\end{enumerate}
\TODO{One sentence of findings, matching the abstract, once the numbers are
final.}

\begin{draftcontrib}{Introduction}
\emph{What is new.} The works we cite each follow one quantity: the distance
the kernel moves \citep{mohamadiWhyYouGrok2024}, the weight norm
\citep{khanhWeightNormSets2026}, the logit scale
\citep{khanhWhatDoesWeight2026}, or the weight-decay rate
\citep{kimGrokkingWeightDecayClock2026,pracherSpectralTheoryGrokking2026}. None
of them separates the growth of the tangent kernel from its rotation. None tests
a scale clock, a rotation clock and a weight-decay clock against each other on
one grid. This report does both, and it adds a placebo test that any claim of a
fixed value at grokking can be given.
\emph{Why it matters.} The 2026 accounts disagree about what times grokking. If
the preliminary findings hold, each account is partly right. Weight decay
shortens the grokking time through the scale of the kernel, width lengthens it
through the rotation, and neither sets the time alone. The findings also qualify
the main evidence for the lazy-to-rich account. At the baseline
\pending{most of the kernel's movement is growth}{N2}, so a large distance does
not by itself show that features were learned. Last, the statistics can be read
from any network during training, without knowing its laziness or its effective
regularisation. Rotation read before any run groks already orders runs by their
grokking time (\pending{concordance 0.71 on held-out cells}{E10}).
\TODO{Search the literature beyond the works cited before calling any of this
the first of its kind.}
\end{draftcontrib}

% =====================================================================
\section{Related Work}\label{sec:related}
\budget{1.2 pages}

\paragraph{Grokking.}
\citet{powerGrokkingGeneralizationOverfitting2022} report grokking on small
algorithmic datasets. \citet{nandaProgressMeasuresGrokking2023} reverse-engineer
the algorithm a small transformer learns for modular addition and split training
into memorisation, circuit formation and clean-up.
\citet{gromovGrokkingModularArithmetic2023} gives closed-form weights that solve
modular addition with a quadratic activation, and observes grokking under
full-batch gradient descent with no explicit regularisation.

\paragraph{Lazy and rich training.}
In the infinite-width limit the tangent kernel stays constant, so the network
trains as a kernel method \citep{jacotNeuralTangentKernel2020}.
\citet{chizatLazyTrainingDifferentiable2020} show that lazy training is a matter
of scale. Multiplying a model's output by a factor $\alpha$ brings its training
closer to that of its linearisation as $\alpha$ grows.
\citet{kumarGrokkingTransitionLazy2024} use this factor to move a network
between the two regimes, and propose that grokking is the transition from lazy
to rich dynamics. \citet{mohamadiWhyYouGrok2024} prove that a
permutation-equivariant kernel method needs a constant fraction of the $p^2$
pairs to learn modular addition, and observe that the kernel moves far from its
initial value only after the training data are fitted.

\paragraph{What sets the grokking time.}
\citet{khanhWeightNormSets2026} find that networks trained with weight decay
grok when the weight norm reaches a critical value, and that clamping the norm
above that value delays grokking exponentially.
\citet{khanhWhatDoesWeight2026} show that under cross-entropy the norm acts
through the scale of the logits, and that under squared error it acts through a
different route. \citet[Cor.~5, Eq.~20, arXiv v1]{kimGrokkingWeightDecayClock2026}
derives, for linear models trained with heavy-ball momentum $\beta$ and weight
decay, a grokking time of about $(1-\beta)/(\eta\lambda)$ times a logarithm of an
initial amplitude, and observes this scaling on modular addition.
\citet{pracherSpectralTheoryGrokking2026} predict, for homogeneous networks
trained on the squared error with weight decay, that the grokking time is set by
the product of the learning rate and the weight decay.

\paragraph{The kernel under weight decay.}
For a network trained by gradient flow with $L_2$ strength $\lambda$, the kernel
evolves as
\begin{align}
  \frac{d\Theta_t(x,x')}{dt} &= -2(k-1)\lambda\,\Theta_t(x,x')
  + T_t(x,x') + T_t(x',x),
  \label{eq:kernelflow}\\
  T_t(x,x') &= -\sum_{x''}\partial_{\mu\nu}f_t(x)\,\partial_\mu f_t(x')\,
  \partial_\nu f_t(x'')\,\ell'(x''),\nonumber
\end{align}
where $k$ is the degree of homogeneity of the network, the sum runs over the
training inputs $x''$ and repeated parameter indices, and $\ell'$ is the
derivative of the loss with respect to the output
\citep[Eqs.~S5 and S6, arXiv v2]{lewkowyczTrainingDynamicsDeep2021}. The weight
decay appears only in the first term, which rescales the kernel. Only $T_t$ can
rotate it. At initialisation $T_0$ has expectation of order $1/N$ in the width
$N$, a result that \citet{lewkowyczTrainingDynamicsDeep2021} take from Dyer and
Gur-Ari (2020). In the infinite-width limit the kernel therefore only decays,
with fixed eigenvectors \citep[Thms.~1 and 2]{lewkowyczTrainingDynamicsDeep2021}.
\TODO{Dyer and Gur-Ari (2020) is named here only as Lewkowycz and Gur-Ari's
source. Add it to ref.bib only after checking it on arXiv.}

\paragraph{Predictions.}
\Cref{tab:predictions} sets out what each account predicts for the three tests
of this report. If rotation sets the clock and $T_t$ stays of order $1/N$ along
training, a wider network rotates more slowly and groks later at a fixed weight
decay. The weight decay should then matter little. A weight-decay clock predicts
the opposite.

\begin{table}[H]
  \centering\small
  \caption{What each account predicts for the three tests. The exponents are
    those of $\log T = c + a\log N + b\log(\eta\lambda)$. ``Not stated'' means
    that the source makes no claim. Entries marked ``ours'' are our reading of
    the source.}
  \label{tab:predictions}
  \begin{tabular}{L{2.6cm}L{3.3cm}L{2.5cm}L{1.6cm}L{2.1cm}}
    \toprule
    Account & Fixed at grokking (Tier 1) & Width exponent $a$ (Tier 2) &
    Decay exponent $b$ (Tier 2) & Early predictor (Tier 3$'$) \\
    \midrule
    Rotation clock (our hypothesis) & $R_t$ reaches a threshold &
    $a>0$, since $T_t$ shrinks with width (ours) & $b\approx0$ & $R_t$ \\
    Weight-decay clock
    \citep{kimGrokkingWeightDecayClock2026,pracherSpectralTheoryGrokking2026} &
    Not stated & $a\approx0$; width can enter Kim's law only inside the
    logarithm (ours) & $b\approx-1$ & $\eta\lambda$ alone \\
    Norm or scale clock
    \citep{khanhWeightNormSets2026,khanhWhatDoesWeight2026} & The weight norm
    reaches a critical value; in our terms $S_t$ at grokking is fixed (ours) &
    Not stated & Not stated & $S_t$ \\
    \bottomrule
  \end{tabular}
\end{table}

\begin{draftcontrib}{Related Work}
\emph{The gap.} \Cref{tab:predictions} shows that the three accounts make
different predictions for the same experiments, yet each was tested on its own.
The accounts that time grokking by weight decay also leave open what sets the
time without it. \Citet[p.~3, arXiv v1]{kimGrokkingWeightDecayClock2026} lists
grokking without explicit weight decay, and the creation of new directions by
feature learning, among the things its mechanism does not explain. Rotation of
the tangent kernel measures one form of that feature learning (our reading).
\emph{Why it matters.} A test that sets the accounts against each other on one
grid shows which part of each one holds.
\end{draftcontrib}

% Cut from the golden report's Related Work, kept in the changelog: spectral
% bias and alignment (Rahaman, Basri, Canatar, Kopitkov, Baratin, Shan), silent
% alignment (Atanasov), the early and late dichotomy (Lyu), Omnigrok (Liu), and
% provable grokking in ridge regression (Xu).

% =====================================================================
\section{Methods}\label{sec:methods}
\budget{1.8 pages}

\subsection{Task, network and training}\label{sec:setup}

The task is modular addition with modulus $p=23$, at the baseline of
\citet[App.~8.3, arXiv v3]{kumarGrokkingTransitionLazy2024}. Each pair $(a,b)$ of
residues is encoded as the concatenation $x_{a,b}\in\mathbb{R}^{2p}$ of the
one-hot codes of $a$ and $b$. Its target $y_{a,b}\in\mathbb{R}^{p}$ is the
one-hot code of $(a+b)\bmod p$. The dataset holds all $p^2=529$ pairs. Each seed
draws its own split, with 476 pairs for training and 53 for testing, and its own
initial weights.

The network is a two-layer perceptron without biases,
\begin{equation}
  f(x;\theta) = N^{-1/2}\,W^{(2)}\phi\bigl(D^{-1/2}W^{(1)}x\bigr),
  \label{eq:network}
\end{equation}
with input dimension $D=2p$, width $N$, weights $W^{(1)}\in\mathbb{R}^{N\times
D}$ and $W^{(2)}\in\mathbb{R}^{p\times N}$ drawn from $\mathcal{N}(0,1)$, and the
rectifier $\phi(h)=\max\{h,0\}$. The factor $N^{-1/2}$ is the NTK
parameterisation of \citet{jacotNeuralTangentKernel2020}. Following
\citet{kumarGrokkingTransitionLazy2024}, the trained predictor is
$\tilde f(x;\theta)=\alpha[f(x;\theta)-f(x;\theta_0)]$ with learning rate
$\eta=\eta_0/\alpha^2$. A larger $\alpha$ gives lazier training. The loss is the
squared error averaged over the outputs and the training pairs. Training is
full-batch gradient descent with weight decay,
$\theta_{k+1}=(1-\eta\lambda)\theta_k-\eta\nabla L(\theta_k)$. For plain
gradient descent this is the same as $L_2$ regularisation
\citep{kimGrokkingWeightDecayClock2026}. \Cref{tab:hyper} lists every setting.

\begin{table}[H]
  \centering\small
  \caption{Every setting of the runs. One setting of $(\alpha,N,\eta\lambda)$ is
    a cell, and each cell has 5 seeds.}
  \label{tab:hyper}
  \begin{tabular}{ll}
    \toprule
    Setting & Value \\
    \midrule
    Modulus, pairs & $p=23$, all 529 pairs; 476 for training and 53 for testing \\
    Split and seeds & Model seed $s\in\{0,\dots,4\}$; the split uses data seed $42+s$ \\
    Network & \Cref{eq:network}, no biases, rectifier, weights from $\mathcal{N}(0,1)$ \\
    Width $N$ & $\{50,100,200,400,800,1600\}$ \\
    Output scale $\alpha$ & $\{0.5,1,2\}$ \\
    Learning rate & $\eta=\eta_0/\alpha^2$ with $\eta_0=100$ \\
    Weight decay $\eta\lambda$ & $\{0,\,3\times10^{-6},\,10^{-5},\,3\times10^{-5},\,10^{-4},\,3\times10^{-4},\,10^{-3},\,3\times10^{-3}\}$ \\
    Optimiser and loss & Full-batch gradient descent, no momentum; mean squared error \\
    Budget & 200\,000 steps \\
    Checkpoints & 307 per run: step 0, 120 log-spaced steps to 200\,000, and every 1\,000 steps \\
    Saved weights & 48 steps per run \\
    Runs and compute & 720 runs; \pending{68\,056 s of wall-clock time}{N1} \\
    \bottomrule
  \end{tabular}
\end{table}
% Checkpoints: dataset_sweep.py takes np.logspace(0, log10(steps), 120) rounded
% to integers (107 distinct steps) together with every 1,000 steps. The manifest gives total_seconds = 68056; confirm whether that
% is wall-clock or summed CPU time, and on which machine.

\subsection{The grid}\label{sec:grid}

Each axis of the grid tests one account. The output scale $\alpha$ sets how lazy
the training is \citep{chizatLazyTrainingDifferentiable2020}. The width $N$ sets
the size of $T_t$ in \eqref{eq:kernelflow}, and so how fast the kernel can
rotate. The weight decay $\eta\lambda$ sets the clock of
\citet{kimGrokkingWeightDecayClock2026}. The grid has $3\times6\times8=144$
cells and 720 runs, and every run finished.

\subsection{Scale, rotation and alignment}\label{sec:statistics}

We measure the tangent kernel of one scalar output, the sum of the logits
divided by $\sqrt p$, $g(x;\theta)=p^{-1/2}\sum_{c}f_c(x;\theta)$. This is the
pseudo-NTK of \citet{mohamadiFastWellFoundedApproximation2022}. On the 53 test
pairs of a seed, the kernel at step $t$ is the matrix with entries
$[K_t]_{ij}=\langle\nabla_\theta g(x_i;\theta_t),\nabla_\theta g(x_j;\theta_t)\rangle$.
The kernel of $\tilde f$ is $\alpha^2$ times the kernel of $f$, so the ratios
below do not depend on which of the two is used. With $Y$ the one-hot targets
of the test pairs and $C=I-\frac1{53}\mathbf{1}\mathbf{1}^\top$, we record
\begin{equation}
  S_t = \frac{\Fnorm{K_t}}{\Fnorm{K_0}},\qquad
  R_t = 1-\frac{\langle K_t,K_0\rangle_F}{\Fnorm{K_t}\,\Fnorm{K_0}},\qquad
  A_t = \frac{\langle CK_tC,\,CYY^\top C\rangle_F}
  {\Fnorm{CK_tC}\,\Fnorm{CYY^\top C}}.
  \label{eq:statistics}
\end{equation}
The scale $S_t$ is 1 at initialisation, and a value above 1 means the kernel
grew. The rotation $R_t$ is the kernel distance of
\citet{fortDeepLearningKernel2020}. It is 0 at initialisation and does not
change when the kernel is rescaled. The alignment $A_t$ is the centred
alignment of \citet{cortesAlgorithmsLearningKernels2024} with the label kernel
$YY^\top$, and it also does not change when the kernel is rescaled.

The usual measure of kernel movement is the relative variation
$D_t=\Fnorm{K_t-K_0}/\Fnorm{K_0}$ \citep{geigerDisentanglingFeatureLazy2020}. It
is a function of the other two,
\begin{equation}
  D_t^2 = S_t^2 - 2S_t(1-R_t) + 1,
  \label{eq:identity}
\end{equation}
(\cref{app:identity}). A kernel that only grows has $R_t=0$ and
$D_t=\lvert S_t-1\rvert$. A large $D_t$ therefore does not show that the kernel
rotated.

\TODO{Add a notation table here if the group wants one: $S_t$, $R_t$, $A_t$,
$T$, $\tstar$, $a$, $b$, with the value of each at initialisation and how to read
it. It costs about a quarter of a page.}

\subsection{Event times}\label{sec:events}

The memorisation step of a run is the first checkpoint at which its training
accuracy reaches 0.99. Its grokking step $T$ is the first checkpoint at which
its test accuracy reaches 0.95. Appendix~\ref{app:levels} repeats the analysis
at the levels 0.80 and 0.90. Both events are counted from step 0. A run that
does not reach a level within 200\,000 steps is kept, and its grokking step is
recorded as censored at the budget. The events are read on the checkpoint grid,
so each one is known to lie between two checkpoints \pending{(at most a
tenth of the event step at $\alpha=1$)}{N3}. We use accuracy rather than loss
because the test loss levels off above zero after the test accuracy is
perfect \pending{(between 0.017 and 0.029 in the baseline cell)}{N4}.

\subsection{Analysis}\label{sec:analysis}

\paragraph{Thresholds (Tier 1).} For each statistic we take its value at the
grokking step of every run that groks, and measure the spread across runs by the
coefficient of variation (CV). A statistic can have a small spread because it
stays in a narrow band all through training. The placebo test separates the two
cases. It reads each run at the grokking step of another run, drawn by a random
permutation with no fixed points, and measures the CV again. The tightening
ratio is the median placebo CV over 2\,000 permutations divided by the CV at the
true events. A ratio near 1 means the statistic is not tied to the event.

\paragraph{Scaling (Tier 2).} We fit the log-normal accelerated failure time
model
\begin{equation}
  \log T = c + a\log N + b\log(\eta\lambda) + \gamma\,\log N\,\log(\eta\lambda)
  + \sigma\varepsilon,\qquad \varepsilon\sim\mathcal{N}(0,1),
  \label{eq:aft}
\end{equation}
by maximum likelihood. A censored run contributes the probability that its
grokking step exceeds the budget
\citep{kalbfleischStatisticalAnalysisFailure2002}. Both logarithms are centred at
the middle of the grid, so $a$ and $b$ are the slopes there. The fits use
lifelines \citep{davidson-pilonLifelinesSurvivalAnalysis2019}. The intervals
resample \TODO{whole cells, and seeds within cells, so that they include the
misfit of the cells about the fitted plane}.

\paragraph{Prediction (Tier 3$'$).} We read $S_t$, $\log R_t$ and $A_t$ at
$\tstar=2\,430$, a saved step before any run reaches the grokking level. Four
model families map them to the grokking step: kernel ridge regression with a
Gaussian kernel, a linear accelerated failure time model with interaction terms,
gradient-boosted trees with the accelerated failure time objective
\citep{chenXGBoostScalableTree2016}, and a gated recurrent unit
\citep{choLearningPhraseRepresentations2014} that reads the statistics at the 25
saved steps up to $\tstar$. Two scores are used. $R^2$ is the share of the
spread in $\log T$ explained on held-out runs. The concordance index is the share
of pairs of runs that a model puts in the right order, so that 0.5 is a coin toss
\citep{harrellEvaluatingYieldMedical1982}. The five seeds of a cell are close
copies, so whole groups are held out: one cell at a time, which tests
interpolation, or one value of $\alpha$ at a time, which tests extrapolation.
Appendix~\ref{app:fitting} gives the settings of each model.

% =====================================================================
\section{Experiments}\label{sec:experiments}

\subsection{Baseline: the reproduction of Kumar et al.}\label{sec:baseline}
\budget{0.8 page}

\textbf{Claim.} At the baseline of \citet{kumarGrokkingTransitionLazy2024} the
network memorises within a few thousand steps and groks much later, and a lazier
network groks later. Over this time the kernel grows several-fold while it
rotates by a few per cent.

\begin{figure}[H]
  \centering
  \placeholderfigure[5.0cm]{E1. Four panels against log step: (a) training and
    test accuracy, (b) $S_t$, (c) $R_t$, (d) $A_t$. Cells with $N=100$,
    $\eta\lambda=0$ and $\alpha\in\{0.5,1,2\}$, one colour per $\alpha$. Median
    over 5 seeds with the range shaded. Hollow marker at memorisation, filled
    marker at grokking.}
  \caption{Training at the baseline of \citet{kumarGrokkingTransitionLazy2024}
    for three values of $\alpha$. \TODO{Final caption.}}
  \label{fig:baseline}
\end{figure}

\TODO{One paragraph. Give the median memorisation and grokking steps for each
$\alpha$, with their ranges and the number of seeds that grok within the budget
(N5). Say that grokking comes later as $\alpha$ grows, as in
\citet{kumarGrokkingTransitionLazy2024}, and point to Appendix~\ref{app:kumar}
for the comparison setting by setting.}

\TODO{One paragraph on the split of $D_t$ by \eqref{eq:identity}. Give the share
of $D_t^2$ carried by the scale at memorisation, at grokking and at the last step
(N2), and the final $S_t$ and $R_t$. This motivates measuring $S_t$ and $R_t$
separately.}

\begin{draftcontrib}{Section 4.1}
\emph{What is new.} The reproduction itself is not new. Its purpose is to
confirm that the setup groks as published, which the proposal set as the gate
for all later work. What is new is the split of the kernel's movement at the
very baseline where the lazy-to-rich account was proposed
\citep{kumarGrokkingTransitionLazy2024}.
\emph{Why it matters.} If \pending{most of $D_t^2$ is scale}{N2}, the distance
that serves as evidence of the transition mostly records growth. A kernel that
only grows keeps its eigenvectors, so its features do not change. Rotation has
to be measured on its own before a large distance can be read as feature
learning.
\end{draftcontrib}

\subsection{Tier 1: is there a threshold?}\label{sec:tier1}
\budget{1.6 pages}

\textbf{Claim.} At a fixed width, grokking happens at a nearly fixed amount of
rotation, and the scale at grokking varies much more. The alignment hardly moves
during training. Across widths, no statistic takes a fixed value at grokking.

\begin{figure}[H]
  \centering
  \placeholderfigure[5.0cm]{E2. Three panels, one for each of $S_t$, $R_t$ and
    $A_t$ at the grokking step. Every run that groks is a point. Left block:
    $N=100$, grouped by $\alpha$ and coloured by $\eta\lambda$. Right block:
    $\alpha=1$, grouped by $N$. Each block is labelled with its CV and
    tightening ratio.}
  \caption{The kernel statistics at the grokking step. \TODO{Final caption.}}
  \label{fig:thresholds}
\end{figure}

\begin{table}[H]
  \centering\small
  \caption{Spread of each statistic across runs that grok, as a coefficient of
    variation, at the grokking step, at the memorisation step and at another
    run's grokking step (placebo). The tightening ratio is the placebo CV over the
    CV at grokking.}
  \label{tab:thresholds}
  \begin{tabular}{llrrrr}
    \toprule
    Runs & Statistic & At grokking & At memorisation & Placebo & Tightening \\
    \midrule
    $N=100$, all $\alpha$ and $\eta\lambda$
      & $S_t$ & \pending{0.48}{E3} & \TODO{} & \pending{0.62}{E3} & \pending{1.3}{E3} \\
      & $R_t$ & \pending{0.20}{E3} & \TODO{} & \pending{0.37}{E3} & \pending{1.9}{E3} \\
      & $A_t$ & \pending{0.08}{E3} & \TODO{} & \pending{0.09}{E3} & \pending{1.1}{E3} \\
    \midrule
    $\alpha=1$, all $N$ and $\eta\lambda$
      & $S_t$ & \pending{0.75}{E3} & \TODO{} & \pending{0.78}{E3} & \pending{1.0}{E3} \\
      & $R_t$ & \pending{0.75}{E3} & \TODO{} & \pending{0.77}{E3} & \pending{1.0}{E3} \\
      & $A_t$ & \pending{0.13}{E3} & \TODO{} & \pending{0.13}{E3} & \pending{1.0}{E3} \\
    \bottomrule
  \end{tabular}
\end{table}
% Planning values: 100 permutations at full checkpoint resolution, 5 Oct 2026.
% E3 uses 2,000 permutations and adds the memorisation column and permutation p.

At a fixed width, $R_t$ is the statistic most tied to the event: its spread at
grokking is \pending{less than half}{E3} that of $S_t$, and it widens
\pending{1.9-fold}{E3} when each run is read at another run's grokking step. The
alignment has a small spread, but the placebo test shows that this comes from
staying in a narrow band. It moves from \pending{0.398}{N6} at initialisation to
\pending{0.406}{N6} at grokking (medians over runs). Across widths no statistic
passes the placebo test. At grokking, $R_t$ falls from \pending{0.108 at $N=50$
to 0.019 at $N=1600$}{N7}, and $S_t$ falls from \pending{4.1 to 0.58}{N7}.
\TODO{Name the cells that reach the $R_t$ band of the grokking runs without
grokking, if any remain on these data (N8).}

Wider networks grok with less rotation but take longer, which suggests that
width sets the rate at which the kernel rotates. \Cref{sec:tier2} tests this.

\begin{draftcontrib}{Tier 1}
\emph{What is new.} This is a threshold test of the tangent kernel at grokking,
with a placebo control, across laziness, width and weight decay. The proposal's
hypothesis, that the alignment crosses a threshold shared by every cell, does not
hold here in that form.
\emph{Why it matters.} Reading each statistic at the true event and at another
run's event separates a value tied to the event from one that stays in a narrow
band. Without that check the alignment would have looked like the best threshold
(\pending{CV 0.08}{E3}). At a fixed width the rotation needed to grok is nearly
fixed, so there the grokking time comes down to how fast the kernel rotates.
Across widths the needed rotation changes, so no single kernel state marks
grokking. Threshold claims of this kind, such as the critical weight norm of
\citet{khanhWeightNormSets2026}, can be given the same check.
\end{draftcontrib}

\subsection{Tier 2: width or weight decay?}\label{sec:tier2}
\budget{2.0 pages}

\textbf{Claim.} The grokking time depends on both width and weight decay, so
neither the rotation clock nor the weight-decay clock of
\cref{tab:predictions} holds on its own. Width acts through the rotation of the
kernel, and weight decay acts through its scale.

\begin{figure}[H]
  \centering
  \placeholderfigure[5.0cm]{E4. Map over $N$ (rows) and $\eta\lambda$ (columns)
    at $\alpha=1$. Colour is the median $\log T$. Each cell is labelled with the
    number of seeds that grok out of 5, and a letter for cells where no seed groks
    (M: memorises only, F: forgets, N: never fits).}
  \caption{Grokking time and outcome over width and weight decay at $\alpha=1$.
    \TODO{Final caption.}}
  \label{fig:map}
\end{figure}

\TODO{One paragraph on \cref{fig:map} (E4). Without weight decay, only $N=50$
and $N=100$ grok within the budget. Inside the range where every seed groks, the
time falls with $\eta\lambda$ at every width. The largest weight decay at which a
width still groks falls as the width grows.}

\begin{figure}[H]
  \centering
  \placeholderfigure[5.0cm]{E5. $\log T$ against $\log(\eta\lambda)$ at
    $\alpha=1$, one line per width, seeds as points, censored runs as open
    markers. Dashed: the fitted plane of \eqref{eq:aft}. Grey reference slopes of
    0 (rotation clock) and $-1$ (weight-decay clock).}
  \caption{How the grokking time scales with weight decay at each width.
    \TODO{Final caption.}}
  \label{fig:scaling}
\end{figure}

\begin{table}[H]
  \centering\small
  \caption{Fits of \eqref{eq:aft} at $\alpha=1$ with their 95\% intervals,
    beside the predictions of \cref{tab:predictions}. The main row uses the
    cells with $10^{-5}\le\eta\lambda\le10^{-4}$, where every seed at every width
    groks. The other rows are checks.}
  \label{tab:scaling}
  \begin{tabular}{lrrrr}
    \toprule
    Fit & $a$ & $b$ & $\gamma$ & $\sigma$ \\
    \midrule
    Rotation clock predicts & $>0$ & $\approx0$ & & \\
    Weight-decay clock predicts & $\approx0$ & $\approx-1$ & & \\
    \midrule
    Main range, level 0.95 & \pending{$+0.24$}{E6} & \pending{$-0.63$}{E6} & \pending{$-0.16$}{E6} & \pending{0.38}{E6} \\
    Adding $\eta\lambda=3\times10^{-4}$ & \pending{$+0.43$}{E6} & \pending{$-0.21$}{E6} & \TODO{} & \TODO{} \\
    Level 0.80 & \pending{$+0.30$}{E6} & \pending{$-0.55$}{E6} & \TODO{} & \TODO{} \\
    Level 0.90 & \pending{$+0.27$}{E6} & \pending{$-0.60$}{E6} & \TODO{} & \TODO{} \\
    Memorisation step & \pending{$-1.53$}{E6} & \TODO{} & \TODO{} & \TODO{} \\
    \bottomrule
  \end{tabular}
\end{table}
% Planning values from repoduced-code/results/tier2_consistency.json, part C.
% Their intervals resample seeds only. E6 refits with intervals that resample
% cells as well, and fills the empty entries.

Both exponents are nonzero. A wider network groks later, as the rotation clock
predicts, and more weight decay shortens the time, as the weight-decay clock
predicts. Neither exponent takes the value its account predicts. Wider networks
memorise sooner, so the width effect on grokking is not a delay in fitting the
training data. \TODO{Interpret the sign of $\gamma$ once E6 is final: a negative
$\gamma$ means that weight decay matters more at large width.}

\begin{figure}[H]
  \centering
  \placeholderfigure[4.0cm]{E7. Two maps over $N$ and $\eta\lambda$ at
    $\alpha=1$: median $\log R_t$ and median $S_t$ at $\tstar$. Merge into
    \cref{fig:map} as extra panels if space is short.}
  \caption{Width sets the rotation of the kernel and weight decay sets its
    scale. \TODO{Final caption.}}
  \label{fig:mechanism}
\end{figure}

\TODO{One paragraph on \cref{fig:mechanism} (E7). Narrow networks rotate the
most, and weight decay shrinks the kernel. Give the share of the spread of
$\log R_t$ and of $S_t$ that each setting explains (N9). This links Tier 2 back
to Tier 1: width sets the rotation and weight decay sets the scale, and the
grokking time depends on both.}

\TODO{One sentence on the intervals: if they hold the cells fixed, say so, and
give the misfit of the cells about the plane beside the spread between seeds.}

\begin{draftcontrib}{Tier 2}
\emph{What is new.} This is a crossed grid of width and weight decay, fitted with
censoring and set against the predictions of a rotation clock and of the
weight-decay clocks of \citet{kimGrokkingWeightDecayClock2026} and
\citet{pracherSpectralTheoryGrokking2026}. The grid includes runs without weight
decay, which \citet[p.~3, arXiv v1]{kimGrokkingWeightDecayClock2026} places
outside the scope of the weight-decay mechanism.
\emph{Why it matters.} If both exponents stay nonzero and neither takes its
predicted value (\cref{tab:scaling}), no single published clock is complete. The
data then point to two routes that act together. Width acts through the
rotation of the kernel, and weight decay acts through its scale
(\cref{fig:mechanism}). This reconciles the accounts, since each one describes a
single route. It also matters in practice. In the NTK parameterisation a wider
network groks later at the same weight decay and
\pending{fails at a smaller weight decay}{E4}, so the weight decay that speeds
grokking has to be chosen together with the width.
\end{draftcontrib}

\subsection{\texorpdfstring{Tier 3$'$}{Tier 3'}: predicting the grokking time from the kernel}\label{sec:tier3}
\budget{1.6 pages}

\textbf{Claim.} Read before any run groks, rotation is the best single predictor
of the grokking time. Scale adds information that rotation lacks, and alignment
adds little. The best models use rotation and scale together and let them
interact.

The settings $\alpha$ and $\eta\lambda$ are the controls of this experiment. In
a general network the laziness and the effective regularisation cannot be read
off directly, while the tangent kernel can be measured. We therefore ask how much
of the grokking time the kernel statistics recover on their own. A model that is
given the settings serves as a reference that knows more than any network would
report.

\begin{table}[H]
  \centering\small
  \caption{Kernel ridge regression on subsets of the statistics at $\tstar$, for
    the runs that grok, scored on held-out cells. The last column is the share of
    the spread of each quantity that lies between cells.}
  \label{tab:subsets}
  \begin{tabular}{lrr@{\qquad}lr}
    \toprule
    Statistics & $R^2$ & Concordance & Quantity & Between cells \\
    \midrule
    $A_t$ & \pending{0.02}{E10} & \pending{0.44}{E10} & $\log T$ & \pending{0.96}{E10} \\
    $S_t$ & \pending{0.27}{E10} & \pending{0.64}{E10} & $S_t$ & \pending{0.99}{E10} \\
    $\log R_t$ & \pending{0.37}{E10} & \pending{0.71}{E10} & $\log R_t$ & \pending{0.99}{E10} \\
    $S_t$ and $\log R_t$ & \pending{0.52}{E10} & \pending{0.76}{E10} & $A_t$ & \pending{0.56}{E10} \\
    $\log R_t$ and $A_t$ & \pending{0.39}{E10} & \pending{0.73}{E10} & & \\
    $S_t$ and $A_t$ & \pending{0.26}{E10} & \pending{0.65}{E10} & & \\
    All three & \pending{0.57}{E10} & \pending{0.78}{E10} & & \\
    \bottomrule
  \end{tabular}
\end{table}
% Planning values: 355 grokked runs with eta lambda <= 1e-3, leave-one-cell-out,
% kernel ridge with fixed penalty 0.1 and width 0.1 (the values notebook 4
% chose). Variance shares from model_fitting/4_nonlinear_exploration.ipynb Part 2.

Taken alone, $\log R_t$ explains the most of the grokking time
(\cref{tab:subsets}). Adding $S_t$ raises $R^2$ from \pending{0.37 to 0.52}{E10},
so rotation does not carry all of the information. Adding $A_t$ to both raises it
only to \pending{0.57}{E10}. The rotation and the scale interact. Growth of the
kernel shortens the grokking time only when the kernel has also rotated
(\cref{fig:surface}; interaction coefficient \pending{$-0.62$, interval $-1.41$
to $-0.12$}{E9}).

\begin{figure}[H]
  \centering
  \placeholderfigure[4.5cm]{E8. Concordance on a held-out $\alpha$ for a ladder of
    models: each statistic alone, the linear model, the linear model with
    interactions, kernel ridge, gradient-boosted trees, and the recurrent network
    on the history of the statistics. A dashed line marks the reference model
    given the settings. Points show the three held-out values of $\alpha$, bars
    their mean.}
  \caption{Predicting the grokking step of runs with an unseen $\alpha$.
    \TODO{Final caption.}}
  \label{fig:ladder}
\end{figure}

\TODO{One paragraph on \cref{fig:ladder} (E8). On an unseen $\alpha$ each
statistic alone ranks runs only a little better than a coin toss
(\pending{0.52 to 0.55}{E8}). The linear model with interactions reaches
\pending{0.62}{E8}, gradient-boosted trees \pending{0.71}{E8} and the recurrent
network on the history \pending{0.77}{E8}. The reference model given the
settings reaches \pending{0.78}{E8}.}

\begin{figure}[H]
  \centering
  \placeholderfigure[4.5cm]{E9. Predicted $\log T$ over $S_t$ and $\log R_t$ at
    $\tstar$, with $A_t$ at its median, from kernel ridge. Runs overlaid as
    points. Shade the region with no runs.}
  \caption{Growth of the kernel shortens the grokking time only when the kernel
    has also rotated. \TODO{Final caption.}}
  \label{fig:surface}
\end{figure}

\begin{draftcontrib}{Tier 3$'$}
\emph{What is new.} The works we cite explain the grokking time by settings,
such as the weight-decay rate, or by quantities they set directly, such as a
clamped weight norm. Here the time is predicted from statistics of the tangent
kernel measured before any run groks. The models are scored on cells and on a
laziness they have not seen, and each statistic is removed in turn.
\emph{Why it matters.} This is a third, separate test of the same answer.
Rotation carries the most information, and scale carries information that
rotation lacks, so neither one is the whole clock. The statistics can be read
from any network without knowing its laziness or its effective regularisation.
Early rotation is therefore a candidate progress measure in the sense of
\citet{nandaProgressMeasuresGrokking2023}: a quantity that moves before the test
accuracy does. Its limits are that the statistics mostly reflect the settings
and do not separate the seeds of one cell (\cref{sec:limitations}).
\end{draftcontrib}

\subsection{\texorpdfstring{Tier 4$'$}{Tier 4'} (stretch, to be decided): theory}\label{sec:tier4}
\budget{0.6 page}

\begin{draftnote}
\TODO{Decide whether this section is worth its 0.6 page, and which theory
direction it takes. Decide after E3, E6 and E7 are final, because the right
direction depends on whether the mixed clock holds. Record the decision in the
changelog, then replace this block with the section.}

\paragraph{Option A: the rate-addition model (the earlier draft).} If a width
route and a weight-decay route act side by side, their rates add,
\begin{equation}
  \frac{1}{T} = r_0 N^{-q} + c\,\eta\lambda ,
  \label{eq:rates}
\end{equation}
and the faster route sets the time. Item E11 fits it with censoring and
compares it with \eqref{eq:aft}.
\emph{Cost:} about a day on existing data.
\emph{Value:} a compact form of the mixed clock that may explain the sign of
$\gamma$.
\emph{Risk:} it rests on three assumptions of our own. They are that $T_t$ stays
of order $1/N$ along training, that the two routes add, and the form $N^{-q}$.
No source derives the model, and a good fit would not prove the mechanism.

\paragraph{Option B: test the premise behind \cref{tab:predictions}.} Two checks
with no free parameters, on existing data (item E17). The first asks whether the
early rate of rotation falls as $1/N$, as the order of $T_0$ suggests. The second
uses the homogeneity of the network. Scaling every weight of \eqref{eq:network}
by $c$ scales its tangent kernel by $c^2$. Weight decay alone, with no gradient,
would therefore give $S_t=(1-\eta\lambda)^{2t}$ and $R_t=0$. This is our
derivation, and it matches the term $-2(k-1)\lambda\Theta_t$ of
\eqref{eq:kernelflow} with $k=2$. Comparing the measured $S_t$ with this curve
shows how much of the change of scale the weight decay accounts for.
\emph{Cost:} low.
\emph{Value:} it tests the premise that options A and C rest on, against a
checked source.
\emph{Risk:} low. It is a check, and it is not a new theory.

\paragraph{Option C: can the rotated kernel alone generalise?} Freeze the
tangent kernel at step $t$ and use it for kernel ridge regression from the
training pairs to the test pairs, with a ridge proportional to the kernel's
trace (item E18). Rescaling the kernel then leaves the predictions unchanged, so
what the frozen kernel can learn depends only on its shape. Suppose the frozen
kernel starts to generalise near the step at which the network groks. Then
rotation sets what can be learned, and scale sets only how fast training gets
there. That would give a mechanism for the mixed clock.
\emph{Cost:} medium. The kernels on all $p^2$ pairs must be rebuilt from the 48
saved weight steps, which lie up to 20\,000 steps apart late in training.
\emph{Value:} the highest of the options, since it ties rotation to
generalisation directly.
\emph{Risk:} medium. A first check is that the frozen initial kernel does not
generalise. If it does, the lazy phase alone could generalise, which would
challenge the lazy-to-rich reading in this setting.

\paragraph{Option D: the closed-form network.} This was the theory tier of the
proposal. It bounds $\lVert T_t\rVert$ for the two-layer network with a quadratic
activation, whose grokked weights are known in closed form
\citep{gromovGrokkingModularArithmetic2023}.
\emph{Cost and risk:} high. It needs new runs with a different activation and new
derivations before 9 November.

\paragraph{Option E: drop the section.} Spend the space on the Tier 2 checks at
$\alpha=0.5$ and $\alpha=2$, or move option B to an appendix.

\paragraph{Suggested order.} Run option B first. It is cheap, and the other
options rest on its premise. If time allows, option C adds the most. Option A is
a short add-on. It summarises the data and does not explain them. Option D is
likely out of reach.

\begin{draftcontrib}{Tier 4$'$, if option C is chosen}
\emph{What is new.} A direct test of whether the rotated tangent kernel alone can
generalise, set against the step at which the network groks.
\emph{Why it matters.} It would explain the mixed clock found in Tiers 1 to
3$'$. The shape of the kernel decides what can be learned, and its scale decides
how fast training gets there.
\end{draftcontrib}
\end{draftnote}

\begin{figure}[H]
  \centering
  \placeholderfigure[4.0cm]{The figure of the chosen option: E11 for option A
    ($1/T$ against $\eta\lambda$ with the fitted lines of \eqref{eq:rates}), E17
    for option B (early rotation rate against $N$, and $S_t$ against the curve
    of weight decay alone), or E18 for option C (test accuracy of the frozen
    kernel against step, with the network's grokking step marked).}
  \caption{\TODO{Caption for the chosen option.}}
  \label{fig:rates}
\end{figure}

% =====================================================================
\section{Limitations}\label{sec:limitations}
\budget{0.6 page}

\TODO{Write each point as one or two sentences.}
\begin{itemize}
  \item Every run uses the NTK parameterisation, one task and one architecture.
    A width effect measured here belongs to this parameterisation.
  \item Each cell has 5 seeds.
  \item The events are known only to the resolution of the checkpoint grid.
  \item The fits in Tier 2 depend on which cells enter, and the decay edge bends
    the plane (\cref{tab:scaling}).
  \item The kernel statistics say nothing about which seed of a cell groks
    first, and the reference model given the settings predicts better than any
    kernel model.
  \item The alignment here uses the kernel of the summed logits. The tangent
    kernel of all $p$ logits may behave differently.
  \item The centred predictor is not homogeneous, so the theorem behind
    \eqref{eq:kernelflow} applies to our network only as a guide.
  \item The proposal also planned a clamped-norm arm, a sweep over $p$ and the
    training fraction, and a cross-entropy arm. They were not run.
  \item The proposal defined the events by the loss. The report uses accuracy
    (\cref{sec:events}).
\end{itemize}

% =====================================================================
\section{Conclusion}\label{sec:conclusion}
\budget{0.3 page}

\TODO{Three to five sentences. Answer the research question: rotation matters
most, but the grokking time is set by rotation and scale together, through width
and weight decay. Restate what is new: the scale and rotation split as timing
variables, the placebo threshold test, the crossed grid, and prediction from
quantities that any network exposes. Name one next step.}

\begin{draftcontrib}{Conclusion}
We asked what sets the grokking time. If the full runs confirm the preliminary
findings, rotation of the tangent kernel matters most, but it is not enough on
its own. Grokking happens at a nearly fixed rotation only within one width. The
grokking time depends on both width and weight decay. Early rotation and scale
together predict it better than either one does alone. Each published clock
therefore captures one route: weight decay acts through the scale of the kernel,
and width acts through its rotation. What is new is the split of the kernel's
movement into scale and rotation as timing variables, the placebo test of
thresholds, the crossed grid of width and weight decay, and prediction from
quantities that any network exposes. The next step is \TODO{the theory
direction chosen in \cref{sec:tier4}, or a second task}.
\end{draftcontrib}

\section*{Contribution}

\TODO{One line for each member, from the group.}
\begin{itemize}
  \item Al-Saffi: \TODO{}
  \item Chong: \TODO{}
  \item D'Mello: \TODO{}
  \item Kienzle: \TODO{}
  \item Rieger: \TODO{}
\end{itemize}

\section*{Acknowledgement}

\TODO{Anyone outside the group who advised on the project.}

\section*{Use of AI tools}

\TODO{The task sheet asks each deliverable to state how AI tools were used. The
group should agree on this text. A draft follows.} The group used an AI assistant
(Claude, by Anthropic) to draft parts of the code, to restructure this report,
and to check references against their arXiv pages. Group members checked every
result and every reference before submission.

% =====================================================================
\appendix

\section{The identity for the variation}\label{app:identity}
\budget{0.5 page (appendices: 5 pages in all)}

Expanding the squared norm gives
$\Fnorm{K_t-K_0}^2=\Fnorm{K_t}^2-2\langle K_t,K_0\rangle_F+\Fnorm{K_0}^2$.
Dividing by $\Fnorm{K_0}^2$ and using
$\langle K_t,K_0\rangle_F=(1-R_t)\Fnorm{K_t}\Fnorm{K_0}$ from
\eqref{eq:statistics} gives $D_t^2=S_t^2-2S_t(1-R_t)+1$. By the Cauchy-Schwarz
inequality, $R_t=0$ exactly when $K_t=cK_0$ for some $c>0$, and then
$D_t=\lvert S_t-1\rvert$. \TODO{Check against the golden report, Appendix A.1.}

\section{The baseline against Kumar et al.}\label{app:kumar}
\budget{0.75 page}

\TODO{A table that sets each setting of the baseline beside
\citet[App.~8.3, arXiv v3]{kumarGrokkingTransitionLazy2024}: task, encoding,
training fraction, width, parameterisation, loss, optimiser, learning rate and
$\alpha$ values. Open the paper and check every entry (CLAUDE.md, Rule 3).
\citet{kumarGrokkingTransitionLazy2024} also sweep $\alpha=1.5$, which our grid
leaves out. Add the grokking steps by $\alpha$ beside their reported values
(E12).}

\section{Outcomes over the whole grid}\label{app:grid}
\budget{1.0 page}

\TODO{One table for each $\alpha$: the seeds that grok out of 5 and the median
grokking step for every cell of $N$ and $\eta\lambda$ (E13).}

\section{A cross-check from a second pipeline}\label{app:edge}
\budget{0.75 page}

\TODO{One paragraph and one figure. A second training pipeline (release
grid-v0.1.0 on branch feat/grid) ran 12 seeds per cell at the baseline with
weight decays near the edge, with a budget of 100\,000 steps. The golden report
gives these results: the grokking time is flat in $\eta\lambda$ below a break
at \pending{$\eta\lambda\approx7.0\times10^{-4}$}{E14}, rises steeply above it,
and no seed groks at \pending{$8.28\times10^{-4}$}{E14}. Export the figure as a
PDF from that pipeline (E14).}

\section{Settings of the prediction models}\label{app:fitting}
\budget{1.0 page}

\TODO{Describe the grouped cross-validation, the search grid of each model, the
architecture and training of the recurrent network, and the random seeds (E15).
Notebooks 3 and 4 in \file{model_fitting/} hold these settings.}

\section{Sensitivity to the grokking level}\label{app:levels}
\budget{0.5 page}

\TODO{Repeat \cref{tab:thresholds} and \cref{tab:scaling} at the levels 0.80 and
0.90 (E16).}

\smaller
\printbibliography

\newpage
Indicate whether all your group members consent for your report to be used as a
teaching resource by including one of the statement below.

\begin{center}
\emph{We give consent for this to be used as a teaching resource.}\\
\emph{We DO NOT give consent for this to be used as a teaching resource.}
\end{center}
\TODO{Keep one of the two statements.}

% =====================================================================
% Draft notes for the group. Delete everything from here to \end{document}
% before submission.
\newpage
\normalsize
\begin{draftnote}
\section*{Draft notes: delete before submission}

\subsection*{1. How the report meets each marking criterion}

The task sheet marks the report on five criteria and asks it to cover six
points. The table gives where this skeleton addresses each one, the mark we
expect if the plan is carried out, and how to raise it.

{\small
\begin{tabular}{L{2.5cm}L{3.9cm}L{2.3cm}L{4.7cm}}
  \toprule
  Criterion & Where it is addressed & Likely mark & How to improve \\
  \midrule
  Clarity and conciseness & One research question. Three statistics on one
  kernel. One regression model and one cross-validation scheme. Each section
  opens with its claim, and the spectral and survival machinery is cut. & Good
  once written. The risk is a heavy load of notation and 8 figures in 12 pages.
  & Add the notation table. Merge \cref{fig:mechanism} into \cref{fig:map}. Use
  one colour for each statistic in every figure. Do a pass for the writing rules
  of CLAUDE.md and for spelling. \\
  Correctness & Censored fits, the placebo test, grouped cross-validation, and
  the caveats stated in \cref{sec:limitations}. & Good, with a risk in Tier 2.
  & Report the check rows of \cref{tab:scaling}. Make the intervals resample
  cells. State the assumptions of the theorem behind \eqref{eq:kernelflow}.
  Label any Tier 4$'$ model that no source derives as the group's own. \\
  Reproducibility & \Cref{tab:hyper}, the grid, the split rule, the compute, and
  one script for each figure, listed in the experiments file. & Good. & Publish a
  small extract of the data. Generate every number as a macro. Add a README that
  maps each figure to its script. \\
  Significance & Three 2026 accounts tested on one grid, with statistics that any
  network exposes. The ``why it matters'' half of each draft contribution block
  (Sections 1, 2, 4.1 to 4.4 and 6) states the case section by section. &
  Medium to good. & Keep the reconciling claim (each clock describes one route)
  only if Tier 2 and E7 confirm it. Say what a practitioner can measure and why
  it helps. Discuss what a single small task can and cannot show. \\
  Novelty & The five contributions listed in the Introduction. The ``what is
  new'' half of each draft contribution block says what each section adds over
  the works cited. & Medium to good. The split into scale and rotation may look
  small. & Search the literature beyond the works cited before claiming any
  first. Contrast the split with the distance of \citet{mohamadiWhyYouGrok2024}
  and with the untested argument of \citet{kumarGrokkingTransitionLazy2024}.
  Use the scope of \citet{kimGrokkingWeightDecayClock2026}, which leaves out
  grokking without weight decay. Repeat what is new in the Conclusion. \\
  Required content & Problem and significance (Sections 1 and 2), steps
  (Section 3), findings (Section 4), limitations (Section 5), conclusion
  (Section 6), contribution, AI use and consent. & Complete. & Check the limits
  of 12 pages and 5 appendix pages, and name the file
  G2\_report\_\textless Name\textgreater.pdf. \\
  \bottomrule
\end{tabular}
}

\subsection*{2. Changelog: decisions and what was cut}

\paragraph{Decisions taken on 5 October 2026.}
\begin{itemize}
  \item All results come from \file{data/}: 720 runs, 5 seeds, 200\,000 steps,
    $\alpha\in\{0.5,1,2\}$, six widths and eight weight decays. The golden report
    used the grid-v0.1.0 databases (12 seeds, 100\,000 steps) instead.
  \item The alignment $A_t$ is the centred alignment of the summed-logit kernel
    on the test pairs (\file{A_sum} in \file{data/}). The golden report on
    main uses the tangent kernel of all $p$ logits, after Baratin et al. The
    statistics $S_t$, $R_t$ and $A_t$ now share one kernel, and the prediction
    notebooks already use this $A_t$.
  \item The main text uses one grokking level, 0.95. The levels 0.80 and 0.90 move
    to Appendix~\ref{app:levels}.
  \item The reproduction of \citet{kumarGrokkingTransitionLazy2024} keeps a short
    section with one figure. The comparison setting by setting moves to
    Appendix~\ref{app:kumar}.
  \item Tier 4$'$ was first drafted as the rate-addition model. It was reopened
    the same day as a decision between five options (\cref{sec:tier4}),
    because that model rests on assumptions of our own.
  \item Each section gained a red draft contribution block that states, as if
    the preliminary findings hold, what is new and why it matters.
  \item The word ``rotation'' is kept for now and will become ``shape'' later.
  \item In Tier 3$'$ the model given the settings is a reference, since a general
    network does not report its laziness or its effective regularisation.
  \item The planning checks changed one claim. The alignment does not pass the
    placebo test, so Tier 1 claims a threshold for rotation only.
\end{itemize}

\paragraph{Cut from the golden report.} Each item can be restored from commit
7da1996 on main, with the command
\file{git show 7da1996:docs/report_golden.tex}.
\begin{itemize}
  \item Section 3.4, Spectral measures: effective rank, trace ratios and the
    target power along eigenvectors. Section 4.3: the cumulative target power
    figure (\file{docs/figures/spectrum.tex}).
  \item Section 3.5, Event times: the Kaplan-Meier estimate with Greenwood's
    interval, the phase labels, the additive model of Aalen and the relative risk
    model with time-dependent covariates. Section 4.3: the table of that model.
    Appendix B, Model checks: the probability plot and the cumulative regression
    functions (\file{docs/figures/probability.tex}, \texttt{aalen.tex}).
  \item Section 3.6, Choice of levels: optimal designs, the censored information
    matrix, the broken line and the sequential design. Section 4.2.2: the scan,
    the design table, the broken-line figure, the survival figure and the
    accelerated failure time table of the weight-decay arm
    (\file{docs/figures/phases.tex}, \texttt{broken.tex},
    \texttt{survival.tex}). A short summary stays in Appendix~\ref{app:edge}.
  \item Section 4.2.1: the $\alpha$ sweep with $\alpha=1.5$ and 12 seeds
    (\file{docs/figures/alpha_sweep.pdf}).
  \item Section 3.2: the first-logit kernel as a second output.
  \item Related Work: the paragraph on kernel alignment and spectral bias, silent
    alignment, the early and late dichotomy, Omnigrok, and provable grokking in
    ridge regression.
  \item The pipeline itself: branch \file{feat/grid}, folder \file{grid/}
    (Snakemake workflow, DuckDB databases, R scripts).
\end{itemize}

\paragraph{Planned in the proposal and not run.} The clamped-norm arm with
$\rho\in[1,2]$, the arm over three primes and three training fractions, the
cross-entropy arm, the statistic $y^\top K^{-1}y$, and events defined by loss
thresholds. The proposal is \texttt{proposal.tex} at the root of the repository.

\paragraph{Superseded analyses.} \file{repoduced-code/tier1_v2.ipynb} and
\file{tier2_width_decay.ipynb} on branch \texttt{jasper} use the old data, with
one split shared by every seed, the trace kernel and test accuracy 1.0. Their
methods are ported, and their numbers are not used.

\paragraph{Open decisions for the group.}
\begin{itemize}
  \item Keep the title, or change it to give the answer.
  \item Keep \cref{fig:mechanism} as its own figure, or merge it into
    \cref{fig:map}.
  \item Whether to add the notation table.
  \item Whether the intervals in Tier 2 should resample cells.
  \item Which theory direction, if any, Tier 4$'$ takes (options A to E in
    \cref{sec:tier4}).
\end{itemize}

\subsection*{3. Reference checks}

Every entry of \texttt{ref.bib} was checked on 5 October 2026 against its arXiv
abstract page, its DOI record on Crossref, or PubMed. Each section and equation
number cited in the text was checked in the PDF of the version named.

{\small
\begin{longtable}{L{3.1cm}L{2.3cm}L{1.9cm}L{6.3cm}}
  \toprule
  Entry & Checked against & Status & Notes \\
  \midrule
  \endhead
  Power et al.\ 2022 & arXiv 2201.02177 & Matches & v1 only. \\
  Nanda et al.\ 2023 & arXiv 2301.05217 & Matches & v3 of 19 Oct 2023 gives the month. \\
  Gromov 2023 & arXiv 2301.02679 & Matches & \\
  Jacot et al.\ 2020 & arXiv 1806.07572 & Matches & Year is v4 (Feb 2020); NeurIPS 2018. \\
  Chizat et al.\ 2020 & arXiv 1812.07956 & Matches & Year is v5 (Jan 2020); NeurIPS 2019. \\
  Kumar et al.\ 2024 & arXiv 2310.06110 & Matches & v3 of 11 Apr 2024. App.~8.3 still to be opened (Appendix~\ref{app:kumar}). \\
  Mohamadi et al.\ 2024 & arXiv 2407.12332 & Matches & ICML 2024. \\
  Mohamadi, Bae, Sutherland 2022 & arXiv 2206.12543 & Matches, entry type changed & The golden bib types it as an article with no journal. It is typed misc here. Published at ICML 2023. \\
  Lewkowycz, Gur-Ari 2021 & arXiv 2006.08643, PDFs of v1 and v2 & Matches & Year is v2 (Jan 2021); NeurIPS 2020. CLAUDE.md gives both 2020 and 2021. The kernel flow is Eq.~S5 and $T_t$ is Eq.~S6 in both v1 and v2. The proposal's draft note, which says they become S8 and S9 in v2, is wrong: S8 in v2 is the infinite-width solution. $E[T_0]=O(n^{-1})$ is credited to Dyer and Gur-Ari (2020), which is not in ref.bib. \\
  Fort et al.\ 2020 & arXiv 2010.15110 & Matches & NeurIPS 2020. \\
  Geiger et al.\ 2020 & Crossref and arXiv 1906.08034 & Matches & JSTAT 2020, 113301. \\
  Cortes et al.\ 2024 & arXiv 1203.0550 & Matches & Year is v3 (Apr 2024). The paper is JMLR 13 (2012), 795--828. \\
  Khanh et al.\ 2026 & arXiv 2606.13753 & Matches arXiv; name order flagged & arXiv lists ``Truong Xuan Khanh'' and parses the surname as Khanh, so ``Khanh et al.''\ matches arXiv. In Vietnamese order the family name is usually written first, which would make it ``Truong et al.''. The co-authors have the same question. The proposal's key is \texttt{truong2026norm}. The group should decide. \\
  Khanh 2026 & arXiv 2606.18465 & Matches arXiv; name order flagged & As above. \\
  Kim 2026 & arXiv 2607.23967 & Matches & Cor.~5 and Eq.~20 checked in the v1 PDF (page 8). The Scope paragraph cited in Section 2 and Tier 2 is on page 3 of v1. It lists grokking without explicit weight decay and the creation of new population-visible directions by feature learning among what the mechanism does not explain. \\
  Pracher et al.\ 2026 & arXiv 2609.26679 & Matches, fields added & The golden bib gives only the DOI. The eprint, month and class were added. v1 of 22 Sep 2026. \\
  Kalbfleisch, Prentice 2002 & Crossref, Open Library & Matches, edition changed & The golden bib gives edition 1. The 2002 Wiley printing is the second edition; the first appeared in 1980. Crossref lists the online version as edition 1, so this needs a check of the book itself. \\
  Davidson-Pilon 2019 & Crossref & Matches & JOSS 4(40), 1317. \\
  Harrell et al.\ 1982 (new) & PubMed 7069920, Crossref & Matches & JAMA 247(18), 2543--2546. Crossref lists only the first author; PubMed lists all five. \\
  Chen, Guestrin 2016 (new) & arXiv 1603.02754 & Matches & KDD 2016. \\
  Cho et al.\ 2014 (new) & arXiv 1406.1078 & Matches & EMNLP 2014. \\
  \bottomrule
\end{longtable}
}
\end{draftnote}

\end{document}
